\section*{Open Questions}
\begin{quote}\small
\noindent\textbf{Status.} \emph{Active live-burden/continuation block.}\par
\noindent\textbf{Block function.} This section separates what is still genuinely open from what survives only as historical labeling, so readers can distinguish live fronts from preserved older checkpoints.
\end{quote}
\addcontentsline{toc}{section}{Open Questions}
%% ============================================================

The structural core, the master equation, and their principal projections are
now explicit. The remaining work no longer concerns the existence of the
framework itself, but its residual quantitative closure, phenomenological
extraction, and spectrum-level completion. To avoid confusion, this section
separates the \emph{current live fronts} from the \emph{historical OP labels}
that are still used in later parts of the text.

\paragraph{Current reading.}
Several historically open OPs should no longer be read as fully open. In
particular, the Koide/triolic block is already structurally closed at the
present level, the Dirac-operator problem on $\Hplusminus$ is no longer open as
an independent operator-existence problem, and the photon-frame
$\Theta$-remainder step of the cosmological block is already closed. The live
questions are therefore the following.

\begin{remark}[Reader precedence for historical release snapshots inside the active line]
The next two synchronized-status blocks reconstruct older public-stage status only for navigation and burden genealogy. They do not overrule the current burden ledger, the current theorematic status, or the current tool/QAM/decoder routing of the active line. Whenever one of those historical snapshots sounds flatter or stronger than the normalized present reading, the normalized present reading governs.
\end{remark}

\begin{remark}[v3.23 synchronized status reading for the optical approximation layer]
\label{rem:v323-synchronized-status-reading-for-the-optical-approximation-layer}
The historical public v3.23.0 release line no longer permits a flat reading in which the new optical sector is either ``already a full optical theory'' or ``still only a loose analogy.'' After the first reduced prism-readout lemma, the residual-diagnostic layer, the carrier-side fit/calibration/validation chain, the reconstruction/correction steps, and the bundled optical approximation package, the optical layer now has its own staged status:
\begin{enumerate}[leftmargin=2em,label=(\roman*)]
  \item a \emph{strong global optical claim} remains open;
  \item a \emph{reduced retained-carrier optical sector} is already available;
  \item a \emph{first-order optical workflow} is already available;
  \item a \emph{bundled first-order optical package} is already available as an active citation surface.
\end{enumerate}
Thus the open-front layer must now distinguish between the still-open stronger global optical theory and the already stabilized first-order optical approximation layer.
\end{remark}

\begin{proposition}[v3.23 optical/open-question status synchronization principle]
\label{prop:v323-optical-open-question-status-synchronization-principle}
At the present theorematic and retained-carrier level, the following synchronized status reading holds for the new optical line:
\begin{enumerate}[leftmargin=2em,label=(\roman*)]
  \item the still-open strong claim is a full global optical theory of the crystal--shadow architecture, including non-reduced anisotropic branches, stronger inversion, and a later convergence theory for higher optical refinements;
  \item what is already available is a reduced prism-readout sector on one retained carrier surface with ordinary wavelength ordering;
  \item what is already available as first-order workflow is the chain consisting of residual diagnostics, carrier-side fit ansatz, spectral calibration, overdetermined validation, effective-index reconstruction, carrier-profile correction, and reduced stability/stopping control;
  \item what is already available as a single active surface is the bundled optical approximation package of Proposition~\ref{prop:first-bundled-optical-approximation-package};
  \item therefore the optical sector should no longer be counted among the fully open foundational fronts of the document, but as a partially stabilized first-order approximation layer with a still-open stronger global continuation.
\end{enumerate}
\end{proposition}

\begin{proof}
Clause (i) records exactly the stronger optical ambitions that the current line has not yet proved: unrestricted anisotropic branching, stronger inversion, and a genuine global convergence theory are still absent. Clause (ii) is supplied by the reduced prism-readout lemma. Clause (iii) is the cumulative content of the residual spectral-diagnostic lemma, the carrier-side fit principle, the calibration principle, the overdetermined consistency principle, the effective-index reconstruction principle, the carrier-profile correction principle, and the reduced stability principle. Clause (iv) is exactly the content of the bundled optical approximation package. Hence the optical line already has a stabilized first-order approximation layer even though the stronger global optical theory remains open.
\end{proof}

\begin{corollary}[Use rule for the synchronized v3.23 optical status reading]
\label{cor:use-rule-for-the-synchronized-v323-optical-status-reading}
Whenever a later argument in the present public v3.23.0 release line needs to classify whether the optical sector is already available for active use, it is sufficient to cite Proposition~\ref{prop:v323-optical-open-question-status-synchronization-principle}. A finer reopening is required only when the argument depends on a stronger global optical theory beyond the reduced retained-carrier approximation layer.
\end{corollary}

\begin{proof}
This is the operational reformulation of Proposition~\ref{prop:v323-optical-open-question-status-synchronization-principle}. The proposition already isolates the precise point at which the optical line ceases to count as merely a loose heuristic while still falling short of a full global optical theory.
\end{proof}


\begin{remark}[Historical v3.22/v3.23 synchronized status reading for the upgraded closure layer]
\label{rem:v322-synchronized-status-reading-upgraded-closure-layer}
The historical public v3.23.0 release line no longer permits a flat reading in which the upgraded closure/interface blocks are simply either ``open'' or ``proved.'' After the public v3.20 theorem-level shadow package, the public v3.21 interface-payoff layer, and the present v3.22 partial-closure steps, their status is now genuinely four-tiered:
\begin{enumerate}[leftmargin=2em,label=(\roman*)]
  \item a \emph{strong deferred claim} remains open;
  \item a \emph{theorem-level shadow} is already available;
  \item an \emph{interface-level operational use} is already available;
  \item for certain carrier-side conclusions, a \emph{restricted partial-closure / lift-back consequence} is already available.
\end{enumerate}
Thus the open-front layer must now distinguish not only between strong and weak forms, but also between theorematic shadow, interface use, and restricted closure consequence.
\end{remark}

\begin{proposition}[v3.22 hypothesis/open-question status synchronization principle]
\label{prop:v322-hypothesis-open-question-status-synchronization-principle}
At the present theorematic and interface level, the following synchronized status reading holds:
\begin{enumerate}[leftmargin=2em,label=(\roman*)]
  \item for Hypothesis~\ref{hyp:hc-meta-generator}, the still-open strong claim is the existence of one explicit compensator-centered generating equation $\mathfrak E_{\mathrm{HC}}=0$, while the theorem-level shadow of a common closure-stable carrier/projection architecture is already available, already usable on the terminal interface, and now participates in the restricted lift-back layer whenever only retained-carrier conclusions are drawn;
  \item for Hypothesis~\ref{hyp:frame-dependent-defect-readout}, the still-open strong claim is the existence of a deeper explicit compensator generator $\mathbf H_C$, while the theorem-level shadow separating visible defect vanishing from carrier extinction is already available, already usable on the terminal interface, and now participates in the no-second-carrier / restricted lift-back layer;
  \item for Hypothesis~\ref{hyp:fluid-readout}, the still-open strong claim is the existence of a concrete coarse-graining map to fluid-type evolution, while the theorem-level shadow demoting fluid-like behavior to a later coarse-grained readout is already available, already usable on the terminal interface, and now participates negatively in the restricted lift-back layer by forbidding replacement of the retained carrier;
  \item for Hypothesis~\ref{hyp:bidirectional-probe-lensing}, the still-open strong claim is a full universal geometric lensing law, while the theorem-level shadow of common-carrier/two-probe coupling is already available, already usable on the terminal interface, and now participates in the no-second-carrier / partial-closure layer;
  \item for the OP~3/OP~5 line, the still-open strong claim is a full global closure theorem, while the historical public v3.23.0 release line already contained more than a weak bridge: it had a first partial closure theorem together with a restricted lift-back to retained carrier-layer structure.
\end{enumerate}
Therefore these upgraded blocks should no longer be counted among the \emph{fully open foundational fronts} of the document. Only their strong deferred forms remain fully open; their weaker theorematic, interface-level, and partial-closure forms are already part of the active line.
\end{proposition}

\begin{proof}
Each clause follows by combining the public v3.20 theorem-level shadows, the public v3.21 interface-level application principles, and where applicable the new v3.22 partial-closure consequences. For the HC meta-field generator hypothesis, this combines the common-carrier/projection shadow and its interface-level use with the present restricted lift-back layer. For frame-dependent defect readout, this combines the defect-shadow theorem, the interface-level defect-face vanishing principle, and the present no-second-carrier/lift-back layer. For fluid-type readout, this combines the fluid-shadow theorem, the interface-level fluid demotion principle, and the present lift-back restriction that forbids carrier replacement. For bidirectional probe lensing, this combines the probe-lensing shadow theorem, the interface-level probe-coupling principle, and the present no-second-carrier/partial-closure layer. Finally, the OP~3/OP~5 line is upgraded from the v3.21 bridge test to the present v3.22 partial closure theorem and restricted lift-back theorem. Hence that historical stage already supported a synchronized four-tier status reading.
\end{proof}

\begin{corollary}[Use rule for the synchronized v3.22 status reading]
\label{cor:use-rule-for-the-synchronized-v322-status-reading}
Whenever a later argument refers back to the historical public v3.23.0 release line and needs to classify whether one of the upgraded closure/interface blocks was still truly open at that stage, it is sufficient to cite Proposition~\ref{prop:v322-hypothesis-open-question-status-synchronization-principle}. A finer reopening is required only when the argument depends on the still-open strong form rather than on the already stabilized theorematic, interface-level, or restricted partial-closure form.
\end{corollary}

\begin{proof}
This is the operational reformulation of Proposition~\ref{prop:v322-hypothesis-open-question-status-synchronization-principle}. The proposition already isolates the precise point at which the upgraded blocks cease to count as fully open foundational fronts.
\end{proof}


\begin{remark}[v3.21 synchronized status reading for the upgraded hypothesis layer]
\label{rem:v321-synchronized-status-reading-upgraded-hypothesis-layer}
The present v\docversion\ release line no longer permits a flat reading in which the upgraded hypothesis blocks are simply either ``fully open'' or ``already proved.'' After the released v3.20 theorem-level shadow package and the historical v3.21 interface payoff layer, their status is now genuinely three-tiered:
\begin{enumerate}[leftmargin=2em,label=(\roman*)]
  \item a \emph{strong deferred claim} remains open;
  \item a \emph{theorem-level shadow} is already available;
  \item an \emph{interface-level operational use} is already available on the terminal surface.
\end{enumerate}
Thus the open-front layer must now distinguish between still-open strong forms and already stabilized weaker forms.
\end{remark}

\begin{proposition}[v3.21 hypothesis/open-question status synchronization principle]
\label{prop:v321-hypothesis-open-question-status-synchronization-principle}
At the present theorematic and interface level, the following synchronized status reading holds:
\begin{enumerate}[leftmargin=2em,label=(\roman*)]
  \item for Hypothesis~\ref{hyp:hc-meta-generator}, the still-open strong claim is the existence of one explicit compensator-centered generating equation $\mathfrak E_{\mathrm{HC}}=0$, while the theorem-level shadow of a common closure-stable carrier/projection architecture is already available and already usable on the terminal interface;
  \item for Hypothesis~\ref{hyp:frame-dependent-defect-readout}, the still-open strong claim is the existence of a deeper explicit compensator generator $\mathbf H_C$, while the theorem-level shadow separating visible defect vanishing from carrier extinction is already available and already usable on the terminal interface;
  \item for Hypothesis~\ref{hyp:fluid-readout}, the still-open strong claim is the existence of a concrete coarse-graining map to fluid-type evolution, while the theorem-level shadow demoting fluid-like behavior to a later coarse-grained readout is already available and already usable on the terminal interface;
  \item for Hypothesis~\ref{hyp:bidirectional-probe-lensing}, the still-open strong claim is a full universal geometric lensing law, while the theorem-level shadow of common-carrier/two-probe coupling is already available and already usable on the terminal interface.
\end{enumerate}
Therefore these four blocks should no longer be counted among the \emph{fully open foundational fronts} of the document. Only their strong deferred forms remain open; their weaker theorematic and interface-operational forms are already part of the active line.
\end{proposition}

\begin{proof}
Each clause follows by combining one theorem-level shadow result with its corresponding interface-level application principle. For Hypothesis~\ref{hyp:hc-meta-generator}, this is the theorem-level shadow of the HC meta-field generator together with the interface-level common-carrier/projection principle. For Hypothesis~\ref{hyp:frame-dependent-defect-readout}, this is the theorem-level shadow of frame-dependent defect readout together with the interface-level defect-face vanishing principle. For Hypothesis~\ref{hyp:fluid-readout}, this is the theorem-level shadow of fluid-type readout together with the interface-level fluid-readout demotion principle. For Hypothesis~\ref{hyp:bidirectional-probe-lensing}, this is the theorem-level shadow of probe-lensing together with the interface-level probe-coupling principle. Hence the four upgraded blocks now share one synchronized three-tier status reading: strong deferred form, theorem-level shadow, and interface-level operational use.
\end{proof}

\begin{corollary}[Use rule for the synchronized v3.21 status reading]
\label{cor:use-rule-for-the-synchronized-v321-status-reading}
Whenever a later argument in the present v\docversion\ release line needs to classify whether one of the four upgraded hypothesis blocks is still truly open, it is sufficient to cite Proposition~\ref{prop:v321-hypothesis-open-question-status-synchronization-principle}. A finer reopening is required only when the argument depends on the still-open strong form rather than on the already stabilized theorematic or interface-level shadow.
\end{corollary}

\begin{proof}
This is the operational reformulation of Proposition~\ref{prop:v321-hypothesis-open-question-status-synchronization-principle}. The proposition already isolates the precise point at which the upgraded blocks cease to count as fully open foundational fronts.
\end{proof}


\begin{enumerate}
  \item \textbf{Residual quantitative completion of the frame scale}
        \textbf{(historical OP~1; also absorbing OP~2(T4) and the former OP~4 coupling).}
        The continuum fixed-point structure is already identified at the
        Hopf/Tribonacci level, and the frame-offset problem is no longer a
        broad undifferentiated open front. What remains open is the explicit
        derivation of the prime-lattice / continuum-limit correction
        $\delta_{\mathrm{CL}}$ (or an equivalent arithmetic correction)
        that transports the continuum candidate to the measured matter-frame
        value.

  \item \textbf{Phenomenology: controlled emergence.}
        The probe--closure field system is now written in explicit
        variational form and closed by the residue constraint
        $\mathfrak R_C=0$.
        Structurally, the preceding sections identify Dirac-, QFT-, and
        Einstein-type regimes as distinguished projections of the same
        closure-stable equation. What remains open is the fully controlled
        derivation of these projections with explicit constants,
        approximation bounds, and quantitative phenomenology.

  \item \textbf{Spectrum: mass hierarchy.}
        The present structural reading is already strong enough to separate the
        qualitative mass-sector logic from the unresolved quantitative
        assignment. The remaining open task is the map from prime/spectral
        index to Standard-Model quantum numbers and measured mass scales.

  \item \textbf{Spectrum: three generations.}
        $S^2\supset S^1\supset S^0$ as three fermion generations is
        structurally supported, but the quantitative generation-splitting law
        is not yet derived.

  \item \textbf{Structural completion: Dixon algebra.}
        A compact reformulation in
        $\mathbb{C}\otimes\HH\otimes\OO$ (64 dimensions) may unify
        the present results into one algebraic object, but this remains a
        later consolidation step rather than a current prerequisite.

  \item \textbf{$\Lambda$-cancellation: reduced photon-frame residual problem}
        \textbf{(historical OP~5, Step~(B) only).}
        The photon-frame formulation and the structural separation
        $\Lambda_{\mathrm{fundamental}}=0$ versus quantitative balance are now
        fixed, and Step~(A) of the $\Theta$--$R$ program is already closed in
        the present text. The remaining open content is Step~(B): the exact
        $\Theta$-sum / $\Theta$--$R$ balance needed to cancel the constant term
        at the photon-frame cutoff.
\end{enumerate}

\paragraph{Historical OP synchronization.}
For the remainder of the document, the historical labels should be read with
these statuses in mind:
\begin{itemize}[leftmargin=2em]
  \item \textbf{Historical OP~3 (Koide / $\theta_K$ block):} closed at the
        present structural level; later work may still densify exposition, but
        it should no longer be listed as a primary open problem.
  \item \textbf{Historical OP~2 (Dirac operator on $\Hplusminus$):} closed as
        an independent operator problem in its existence, self-adjointness,
        and GOE-class parts; only the former frame-offset part survives, and
        it is absorbed into Item~1 above.
  \item \textbf{Historical OP~5:} reduced to Step~(B) only; Step~(A) is no
        longer open.
\end{itemize}




\subsection*{Phase IV-I: Physical Unfolding, Observer Frames, and Interface Regimes}
\addcontentsline{toc}{subsection}{Phase IV-I: Physical Unfolding, Observer Frames, and Interface Regimes}

\begin{lemma}[Observer-frame lemma]
\label{lem:observer-frame-lemma}
Within the normalized structural architecture, an observer frame is not introduced as an external background datum. Rather, it is the stabilized readout frame selected by the combined action of the cascade, the closure filter, and the admissibility-preserving transport across the grouped layers. In particular, any observer description compatible with the kernel must arise as a frame-shifted but ground-linked reading of the same underlying body.
\end{lemma}

\begin{proof}
The semantic kernel identifies all admissible descriptions as readings of a single underlying structure. Frame-shifted observation therefore does not create an ontologically new sector. By the forcing chain established in Phases~I--III, any admissible readout is first re-identified by the compensator, then filtered into a closure-stable class, then transported through the grouped architecture without ontology drift. The observer frame is therefore fixed only at the level of admissible readout, not posited independently of the structure.
\end{proof}

\begin{lemma}[Mass/spectral reading lemma]
\label{lem:mass-spectral-reading-lemma}
Any admissible mass or hierarchy reading in the Companion must be interpreted first as a spectral or sectoral readout of the normalized structural kernel, and only secondarily as a phenomenological mass assignment. In particular, the transition from spectral organization to mass hierarchy is a realization map, not a new generative axiom.
\end{lemma}

\begin{proof}
The generative layer produces admissible sectors and their closure relations, while the projection/translation layer supplies readout rules and collapse-fixed identifications. Therefore any mass reading that appears in the field/probe layer inherits its legitimacy from a prior spectral or sectoral differentiation. This excludes the reverse move of inserting an independent mass ontology into the generative layer. The mass interpretation is thus subordinate to the spectral structure and must be treated as a downstream realization.
\end{proof}

\begin{lemma}[Generation-splitting lemma]
\label{lem:generation-splitting-lemma}
If the triolic and grouped-layer architecture is preserved across admissible projections and field realizations, then any generation splitting must appear as a structured decomposition of an already admissible spectral body rather than as a free multiplication of unrelated sectors.
\end{lemma}

\begin{proof}
Triolic minimality and orthogonal anchoring prevent arbitrary proliferation of sectors at the generative level. The transport rules for $\mathcal M$ preserve admissibility and closure, and the no-ontology-drift lemma excludes the creation of new bodies during layer transfer. Hence any multiplicity such as a generation-like splitting must be read as an organized decomposition of the inherited spectral structure. It is therefore constrained by the prior closure and transport logic.
\end{proof}

\begin{theorem}[Standard-model interface theorem]
\label{thm:standard-model-interface-theorem}
The normalized Companion architecture defines a one-way interface to Standard-Model-style phenomenology: admissible structures may be projected into observer, spectral, and field realizations compatible with particle-physics language, but the Companion introduces no independent Standard-Model ontology at the generative level. Consequently, the appropriate first target of the unfolding is not a complete reconstruction of the Standard Model, but a structurally controlled interface map from the grouped kernel into physically readable sectors.
\end{theorem}

\begin{proof}
By the observer-frame, mass/spectral, and generation-splitting lemmas, all physically readable sectors arise only after the generative core has been passed through closure, confinement, cascade readout, tensor realization, and grouped transport. The field/probe layer therefore supplies an interface regime in which known physical language may be attached to admissible sectors, but the kernel itself remains prior to that language. The Standard Model can thus enter only through an interface map respecting admissibility, closure, and no-ontology drift. This proves that the correct unfolding target is an interface theorem rather than a premature full identification.
\end{proof}

\begin{corollary}[First physical unfolding chain]
\label{cor:first-physical-unfolding-chain}
The first physical unfolding chain of the Companion is
\[
\begin{aligned}
&\text{observer-frame emergence}
  \Longrightarrow \text{spectral/mass readout}\\
&\text{generation-compatible splitting}
  \Longrightarrow \text{controlled Standard-Model interface}.
\end{aligned}
\]
Thus the physical interpretation layer is admitted only after the structural, closure, tensorial, and grouped-layer chains have already been stabilized.
\end{corollary}


\subsection*{Phase IV-I$'$: Refined Physical Interface and Realization Control}
\addcontentsline{toc}{subsection}{Phase IV-I': Refined Physical Interface and Realization Control}

\begin{lemma}[Observer covariance under admissible transport]
\label{lem:observer-covariance-admissible-transport}
If two observer descriptions are related by admissibility-preserving transport through the grouped architecture, then they represent frame-shifted realizations of the same underlying body and cannot define distinct generative ontologies.
\end{lemma}

\begin{proof}
The no-ontology-drift result of Phase~II-F implies that transport across $\mathfrak G$, $\mathfrak P$, and $\mathfrak F$ preserves the underlying structural identity. By the observer-frame lemma, any observer description is already a frame-shifted readout of the same admissible kernel. Therefore two observer descriptions related by admissible transport differ only at the level of readout covariance and not at the level of generative content.
\end{proof}

\begin{proposition}[Spectral-to-mass realization map]
\label{prop:spectral-to-mass-realization-map}
There exists a structurally controlled realization map
\[
\mathcal R_{m}: \mathcal S_{\mathrm{spec}} \longrightarrow \mathcal S_{\mathrm{mass}}
\]
from admissible spectral sectors to admissible mass-readout sectors, where $\mathcal R_{m}$ is downstream of closure, cascade readout, and grouped transport.
\end{proposition}

\begin{proof}
By the mass/spectral reading lemma, any mass assignment must be subordinated to a prior spectral organization. The admissible spectral sectors are produced only after the closure-stable body has been transported through the projection/translation layer and read out through the cascade. Hence the passage to a mass language is not primitive but realized by a downstream map from already admissible spectral sectors. Denote this realization map by $\mathcal R_m$. Its domain is therefore restricted to $\mathcal S_{\mathrm{spec}}$ and its codomain inherits admissibility from the upstream structural chain.
\end{proof}

\begin{proposition}[Generation-compatible decomposition]
\label{prop:generation-compatible-decomposition}
If a generation-like multiplicity appears in the physical layer, then it must be expressible as an admissible decomposition of the image of $\mathcal R_m$ rather than as an unconstrained replication of primitive sectors.
\end{proposition}

\begin{proof}
The generation-splitting lemma already excludes free multiplication of unrelated sectors. Since the physical layer is reached only after grouped transport and spectral-to-mass realization, any multiplicity that survives into the interface regime must be inherited from the already admissible image of $\mathcal R_m$. Thus a generation-like structure can only be accepted when it is realizable as a constrained decomposition of an upstream admissible spectral body.
\end{proof}

\begin{theorem}[Controlled physical interface functor]
\label{thm:controlled-physical-interface-functor}
The first physically readable regime of the Companion is determined by a controlled interface functor
\[
\mathfrak I_{\mathrm{phys}}:
(\mathfrak G,\mathfrak P,\mathfrak F)
\longrightarrow
\mathcal O_{\mathrm{phys}},
\]
which assigns observer, spectral, mass, and generation-level readouts to admissible structural sectors while preserving closure compatibility and ontology non-drift.
\end{theorem}

\begin{proof}
The observer covariance lemma identifies physically readable frames as transport-compatible readouts of the same body. The spectral-to-mass realization proposition supplies a downstream realization map from spectral sectors to mass sectors, while the generation-compatible decomposition proposition constrains multiplicity to inherited admissible decompositions. Collecting these assignments yields a functorial interface from the grouped structural architecture to a space of physically readable observables $\mathcal O_{\mathrm{phys}}$. Since every step is subordinated to admissibility, closure, and no-ontology drift, the resulting interface is controlled rather than ontologically generative.
\end{proof}

\begin{corollary}[Refined physical interface chain]
\label{cor:refined-physical-interface-chain}
The refined physical interface chain is
\[
\begin{aligned}
&\text{observer covariance}
  \Longrightarrow \text{spectral-to-mass realization}\\
&\text{generation-compatible decomposition}
  \Longrightarrow \text{controlled interface functor}.
\end{aligned}
\]
Hence the physical language of the Companion is introduced only as a regulated readout of the normalized structural architecture.
\end{corollary}

%% ============================================================
%% APPENDIX A — Formal Proof Supplement
%% ============================================================







\subsection*{Phase IV-J: Spectral Readout, Trace Reconstruction, and Observable Selection}
\addcontentsline{toc}{subsection}{Phase IV-J: Spectral Readout, Trace Reconstruction, and Observable Selection}

\begin{lemma}[Spectral admissibility inheritance]
\label{lem:spectral-admissibility-inheritance}
If a physically readable sector is admitted by the controlled interface functor, then its spectral representative must inherit admissibility from the upstream closure-stable and transport-compatible structural body. In particular, no spectral mode may enter the readout regime unless it descends from an already admissible structural sector.
\end{lemma}

\begin{proof}
By the semantic--tensor interface and the grouped transport chain, admissibility is preserved from the generative layer through cascade readout and into the physically readable interface regime. The controlled physical interface functor therefore cannot attach an observable meaning to a mode that was not already admitted upstream. Hence spectral readability inherits admissibility rather than creating it locally.
\end{proof}

\begin{proposition}[Trace-compatible spectral decomposition]
\label{prop:trace-compatible-spectral-decomposition}
Every admissible spectral body in the readout regime admits a trace-compatible decomposition into modes whose contribution respects the closure filter, the grouped transport constraints, and the observer-covariant interface conditions.
\end{proposition}

\begin{proof}
The historical spectral and trace cores already isolate the need for admissible mode selection and trace-compatible decomposition. In the normalized architecture, these requirements are now downstream of closure, cascade, and grouped transport. Therefore any accepted spectral decomposition must preserve the admissible body identified by the upstream chain and cannot separate into trace-incompatible or closure-breaking components. This yields a trace-compatible mode decomposition of every admissible spectral body.
\end{proof}

\begin{proposition}[Observable extraction map]
\label{prop:observable-extraction-map}
There exists a structurally controlled observable extraction map
\[
\mathcal R_{\mathrm{obs}}: \mathcal S_{\mathrm{spec}}^{\mathrm{adm}} \longrightarrow \mathcal O_{\mathrm{phys}}^{\mathrm{adm}},
\]
which assigns observable readouts only to trace-compatible admissible spectral sectors.
\end{proposition}

\begin{proof}
Phase~IV-I$'$ introduced a controlled interface functor from the grouped architecture into physically readable observables. Restricting that interface to the spectrally admissible, trace-compatible body provided by Proposition~\ref{prop:trace-compatible-spectral-decomposition} yields a downstream extraction map from admissible spectral sectors to admissible observables. Since its domain is already filtered by closure, transport, and trace compatibility, the map does not generate new observables freely but only extracts structurally sanctioned ones.
\end{proof}

\begin{theorem}[Mode-selection and observable consistency theorem]
\label{thm:mode-selection-observable-consistency}
Mode selection in the Companion is admissible only when it is simultaneously spectral, trace-compatible, closure-preserving, and observer-covariant. Under these conditions, the observable extraction map produces a consistent readout regime in which selected modes can be interpreted without violating ontology non-drift.
\end{theorem}

\begin{proof}
Spectral admissibility inheritance excludes modes that do not descend from the admissible structural chain. Trace-compatible spectral decomposition excludes decompositions that would destroy the integrity of the admissible spectral body. The controlled interface functor and observer covariance from Phase~IV-I$'$ ensure that the resulting readout remains a covariance of one and the same body rather than a proliferation of ontologies. Therefore mode selection is consistent exactly when all four constraints are met, and in that regime observable extraction remains structurally controlled.
\end{proof}

\begin{corollary}[Spectral readout reconstruction chain]
\label{cor:spectral-readout-reconstruction-chain}
The next spectral unfolding chain is
\[
\begin{aligned}
&\text{spectral admissibility inheritance}
\Longrightarrow
\text{trace-compatible decomposition}\\
&\Longrightarrow
\text{observable extraction}
\Longrightarrow
\text{consistent mode-selected readout}.
\end{aligned}
\]
Thus the spectral and trace layer of the Companion is admitted only after closure, grouped transport, and controlled physical interfacing have already been stabilized.
\end{corollary}


% -- relocated late main-text phases IV-K to IV-O --


\subsection*{Phase IV-K: Resolvent Gatekeeping and Operator-Controlled Spectral Access}
\addcontentsline{toc}{subsection}{Phase IV-K: Resolvent Gatekeeping and Operator-Controlled Spectral Access}

\begin{lemma}[Resolvent admissibility lemma]\label{lem:resolvent_admissibility}
Let $\mathcal R(\lambda)=(\lambda\mathbf 1-\mathcal A)^{-1}$ denote a formal resolvent on the admissible spectral layer. If the underlying mode family is closure-stable and trace-compatible, then resolvent access is admissible only on those channels for which the induced transport preserves the grouped-layer constraints and does not reintroduce ontology drift. In particular, admissible resolvent access is a restricted rather than global spectral operation.
\end{lemma}
\begin{proof}
The semantic kernel already constrains admissible transport by closure preservation, grouped transport, and controlled readout. A resolvent therefore cannot be treated as an arbitrary inversion device: it is allowed only where the corresponding operator action remains inside the admissible closure class and respects the readout filter established in the preceding phases. Hence admissibility is inherited by resolvent access only on the structurally licensed channels.
\end{proof}

\begin{proposition}[Gatekeeping proposition for admissible modes]\label{prop:resolvent_gatekeeping}
The resolvent layer acts as a gatekeeping mechanism: admissible modes are precisely those spectral branches for which resolvent access is compatible with closure, trace reconstruction, and grouped-layer transport. Non-admissible branches are excluded before observable extraction.
\end{proposition}
\begin{proof}
By Lemma~\ref{lem:resolvent_admissibility}, resolvent access is already restricted to structure-preserving channels. Combined with the trace-compatible mode selection of Phase~IV-J, this yields a gatekeeping principle: only those modes that remain compatible with closure and transport can pass into the observable sector. The exclusion of non-admissible branches is therefore structural rather than ad hoc.
\end{proof}

\begin{theorem}[Operator-controlled trace access theorem]\label{thm:operator_controlled_trace_access}
Assume that the spectral sector is admissible, the selected modes are resolvent-gated, and the trace reconstruction remains compatible with the semantic--tensor interface. Then trace access is operator-controlled: every observable trace contribution arises through an admissible operator channel and not through uncontrolled spectral leakage.
\end{theorem}
\begin{proof}
The preceding proposition ensures that only admissible modes are available for readout. The semantic--tensor interface then identifies the operator realization of the same admissible content, while closure preservation prevents leakage into non-licensed channels. Consequently, the trace contributions that survive observable extraction are exactly those transmitted through operator-controlled admissible channels.
\end{proof}

\begin{corollary}[Stable observable channel corollary]\label{cor:stable_observable_channel}
Under admissible resolvent gatekeeping, the observable readout sector forms a stable channel: spectral access, trace reconstruction, and mode selection remain synchronized under the same structural admissibility constraints.
\end{corollary}
\begin{proof}
This follows directly from Theorem~\ref{thm:operator_controlled_trace_access}: once admissible access, admissible mode selection, and admissible trace transport coincide, the resulting observable sector is stable under the same structural constraints.
\end{proof}

\begin{remark}[Resolvent-to-readout chain]\label{rem:resolvent_to_readout_chain}
The operator-controlled spectral access chain may therefore be summarized as
\[
\begin{aligned}
\text{resolvent admissibility}
&\Longrightarrow \text{gatekept mode access}\\
&\Longrightarrow \text{operator-controlled trace access}\\
&\Longrightarrow \text{stable observable channel}.
\end{aligned}
\]
This phase extends the preceding spectral readout layer by showing that observable extraction is not only spectrally filtered but also operatorically regulated.
\end{remark}



\subsection*{Phase IV-L: Semigroup Stability and Lyapunov-Controlled Readout Persistence}
\addcontentsline{toc}{subsection}{Phase IV-L: Semigroup Stability and Lyapunov-Controlled Readout Persistence}

\begin{lemma}[Semigroup admissibility lemma]\label{lem:semigroup_admissibility}
Let $(T_t)_{t\ge 0}$ denote a formal semigroup acting on the admissible readout sector. If the initial spectral channel is closure-stable, resolvent-gated, and trace-compatible, then admissible semigroup evolution is restricted to those trajectories for which grouped transport, closure preservation, and observable compatibility remain jointly valid for all admissible times.
\end{lemma}
\begin{proof}
The preceding phases identify the observable channel only after admissibility, gatekeeping, and operator-controlled trace access have been enforced. A semigroup acting on this sector may therefore propagate states only if the induced time evolution respects the same admissibility filter. Otherwise the evolution would exit the closure-preserving and trace-compatible regime and would no longer represent a structurally licensed readout trajectory. Hence admissible semigroup evolution is constrained to the channels that preserve the previously established structural conditions for all admissible times.
\end{proof}

\begin{proposition}[Lyapunov gate proposition]\label{prop:lyapunov_gate}
Assume that an admissible Lyapunov functional $\mathcal L$ exists on the gated readout sector. If $\mathcal L$ is nonincreasing along the admissible semigroup trajectories, then unstable leakage into non-admissible channels is structurally suppressed, and the gated observable channel remains dynamically protected.
\end{proposition}
\begin{proof}
A nonincreasing Lyapunov functional excludes the spontaneous amplification of perturbations that would carry the evolution outside the admissible channel. Since admissibility already forbids ontology drift, closure violation, and non-gated spectral leakage, Lyapunov monotonicity upgrades these static constraints to a dynamical protection principle. The readout channel is therefore not merely selected once but remains protected under admissible evolution.
\end{proof}

\begin{theorem}[Persistent observable channel theorem]\label{thm:persistent_observable_channel}
Suppose that the readout sector is resolvent-gated, operator-controlled, and semigroup-admissible, and that a Lyapunov functional controls admissible perturbations. Then the observable channel is persistent: admissible mode selection, trace reconstruction, and grouped transport remain synchronized along the admissible evolution, and the readout sector stays dynamically stable against structural leakage.
\end{theorem}
\begin{proof}
By Lemma~\ref{lem:semigroup_admissibility}, admissible semigroup evolution preserves the structurally licensed readout sector. Proposition~\ref{prop:lyapunov_gate} then prevents the admissible trajectories from developing unstable excursions into non-gated channels. Combined with the operator-controlled access theorem and the stable observable channel corollary from Phase~IV-K, this yields persistence of the observable channel: the same admissibility constraints that select the readout sector also stabilize it under admissible evolution.
\end{proof}

\begin{corollary}[Long-time readout stability corollary]\label{cor:long_time_readout_stability}
Under admissible semigroup evolution and Lyapunov control, the observable readout channel remains stable at long admissible times; the selected spectral and trace-compatible sector is therefore not merely accessible but persistently readable.
\end{corollary}
\begin{proof}
This follows directly from Theorem~\ref{thm:persistent_observable_channel}: persistence of the admissible channel implies long-time stability of the same gated and closure-compatible readout sector.
\end{proof}

\begin{remark}[Semigroup-to-stability chain]\label{rem:semigroup_to_stability_chain}
The late-stage stability chain may therefore be summarized as
\[
\begin{aligned}
\text{semigroup admissibility}
&\Longrightarrow \text{Lyapunov-controlled gated evolution}\\
&\Longrightarrow \text{persistent observable channel}\\
&\Longrightarrow \text{long-time readout stability}.
\end{aligned}
\]
This phase completes the operatoric readout picture by showing that admissible channels are not only spectrally and resolvent-gated, but also dynamically stabilized over admissible evolution scales.
\end{remark}



\subsection*{Phase IV-M: Channel Fusion and Multi-Readout Consistency}
\addcontentsline{toc}{subsection}{Phase IV-M: Channel Fusion and Multi-Readout Consistency}

\begin{lemma}[Admissible channel fusion lemma]\label{lem:admissible_channel_fusion}
Suppose that two or more observable channels are individually admissible, resolvent-gated, and semigroup-stable. Then a fused readout may be admitted only if the fusion procedure preserves the common closure filter, the grouped transport discipline, and the no-ontology-drift condition. In particular, fusion is admissible only as a compatibility-preserving synthesis of already licensed channels.
\end{lemma}
\begin{proof}
The previous phases guarantee admissibility, operatoric access, and dynamical persistence for each channel separately. A fusion mechanism cannot be treated as an unrestricted superposition principle, since otherwise it could reintroduce hidden drift between channels that were only separately controlled. Therefore fusion is admissible only when the same closure, transport, and ontology constraints remain valid for the synthesized channel. In that case fusion does not create a new structural body, but re-expresses a common admissible body through multiple synchronized readouts.
\end{proof}

\begin{proposition}[Cross-channel consistency proposition]\label{prop:cross_channel_consistency}
If admissible channels are fused through a compatibility-preserving synthesis, then their spectral, trace-based, and long-time readout contents must agree on the shared admissible sector. Any discrepancy beyond the admitted transport freedom signals a non-admissible fusion rather than a genuine multi-readout effect.
\end{proposition}
\begin{proof}
Each admissible channel already carries its observable content through the same semantic--tensor interface, the same resolvent gatekeeping, and the same semigroup-stable evolution class. Hence a compatibility-preserving fusion may only combine channels whose readout contents coincide on the overlap of their admissible sectors. If an apparent fusion produces incompatible observable contents beyond the allowed transport freedom, then the synthesis has violated the common admissibility constraints and is therefore not structurally licensed.
\end{proof}

\begin{theorem}[Multi-readout consistency theorem]\label{thm:multi_readout_consistency}
Under admissible channel fusion, the Companion admits a controlled multi-readout regime: distinct readout channels may coexist, but only as synchronized presentations of one and the same admissible structural body. Consequently, multi-readout consistency is equivalent to fusion without ontology drift.
\end{theorem}
\begin{proof}
By Lemma~\ref{lem:admissible_channel_fusion}, channel fusion is allowed only as a synthesis that preserves closure, grouped transport, and non-drift. Proposition~\ref{prop:cross_channel_consistency} then forces agreement of spectral, trace, and persistent readout content on the common admissible sector. Therefore multiple channels can coexist only as mutually compatible renderings of the same admissible body. The controlled multi-readout regime is thus precisely the regime in which fusion introduces no ontology drift.
\end{proof}

\begin{corollary}[Fusion-to-consistency chain]\label{cor:fusion_to_consistency_chain}
The late-stage fusion chain may therefore be summarized as
\[
\begin{aligned}
\text{admissible channel fusion}
&\Longrightarrow \text{cross-channel consistency}\\
&\Longrightarrow \text{controlled multi-readout regime}\\
&\Longrightarrow \text{fusion without ontology drift}.
\end{aligned}
\]
This phase closes the late readout unfolding by showing that stable observable channels may be combined only through a synthesis that preserves the unique admissible structural body across all readout presentations.
\end{corollary}


\subsection*{Phase IV-N: Unified Observable Readout Synthesis}
\addcontentsline{toc}{subsection}{Phase IV-N: Unified Observable Readout Synthesis}

\begin{proposition}[Unified observable readout synthesis]
\label{prop:unified-observable-readout}
Let the late readout regime be governed by the spectral admissibility, resolvent gatekeeping, semigroup persistence, and admissible channel-fusion principles established in Phases IV-J through IV-M. Then the resulting observable sector admits a unified synthesis map in which spectral access, operator control, long-time persistence, and multi-readout consistency are read as mutually compatible aspects of one and the same admissible observable architecture.
\end{proposition}

\begin{proof}
Phase IV-J supplies the spectral and trace-compatible decomposition of admissible modes. Phase IV-K restricts this decomposition to resolvent-accessible channels only, so that observable extraction is not performed on arbitrary modes but only on structurally permitted ones. Phase IV-L then upgrades admissible access to persistence under semigroup evolution and Lyapunov control, which turns transient access into stable readout inheritance over time. Finally, Phase IV-M shows that multiple such persistent channels can be fused without inducing cross-channel inconsistency or ontological drift. Therefore the late readout regime is not a mere juxtaposition of separate spectral, operator, dynamical, and fusion statements; it forms a single admissible synthesis in which observable extraction is structurally unified.\qedhere
\end{proof}

\begin{corollary}[Late readout master chain]
The late readout architecture admits the compressed chain
\[
\begin{aligned}
&\text{spectral admissibility}
\Longrightarrow \text{resolvent-controlled access}\\
&\Longrightarrow \text{semigroup/Lyapunov persistence}
\Longrightarrow \text{multi-readout consistency}\\
&\Longrightarrow \text{unified observable synthesis.}
\end{aligned}
\]
\end{corollary}

\begin{remark}[Interpretive role of the late synthesis]
This synthesis does not introduce a new ontology behind the observable layer. Rather, it states that the different late readout mechanisms are to be interpreted as coordinated access constraints and persistence conditions on one admissible observable sector generated by the normalized structural core.
\end{remark}



\subsection*{Phase IV-O: Return Bridge to the Generative Core and the Structural Master Equation}
\addcontentsline{toc}{subsection}{Phase IV-O: Return Bridge to the Generative Core and the Structural Master Equation}

\begin{proposition}[Return bridge from unified readout to the generative core]\label{prop:return_bridge_unified_readout}
The unified observable readout synthesized in Phase IV-N introduces no new generative assumption. Rather, it is a late readout layer of the same generative structure determined by the one-body condition, triolic stabilization, orthogonal anchoring, frame-shifted observation, and the HC--closure--confinement chain.
\end{proposition}

\begin{proof}
The earlier unfolding phases establish that admissible readout channels arise only after re-identification, compensator forcing, closure stabilization, and higher confinement. Phases IV-J through IV-N then refine how admissible observable content is accessed, filtered, stabilized, fused, and synthesized. None of those later steps introduces a new generating source; they remain constrained by the same admissibility conditions inherited from the generative layer. Hence the late observable synthesis is a return readout of the original structure rather than an independent origin.
\end{proof}

\begin{theorem}[Master-equation readout closure]\label{thm:master_equation_readout_closure}
Let the structural master equation be the compressed dynamical expression of the admissible generative core. Then the late readout chain
\[
\begin{aligned}
&\text{spectral admissibility}
\Longrightarrow \text{resolvent-controlled access}
\\
&\Longrightarrow \text{semigroup/Lyapunov persistence}
\Longrightarrow \text{multi-readout consistency}
\\
&\Longrightarrow \text{unified observable synthesis}
\end{aligned}
\]
closes back into the same master-equation-governed structure, in the sense that every admissible late observable synthesis remains a governed readout of the original master-equation sector.
\end{theorem}

\begin{proof}
By the semantic/tensor interface and grouped transport phases, admissibility is preserved under the Millennium transport between the generative, projection, and realization layers. The readout phases IV-J through IV-N add only progressively refined access conditions: spectral admissibility, resolvent gatekeeping, semigroup persistence, and fusion consistency. Since each step is formulated as structure-preserving and closure-compatible, the final observable synthesis remains within the same admissible sector selected by the master equation. Therefore the readout chain closes into the same governing structure and should be interpreted as its late observable face.
\end{proof}

\begin{remark}[Interpretive closure of the unfolding program]
The point of the return bridge is to prevent the late readout machinery from being mistaken for a detached auxiliary theory or a second source. The late phases do not compete with the generative core; they complete its observable unfolding. In this sense the overall program closes as
\[
\begin{aligned}
&\text{generative core}
\Longrightarrow \text{HC/closure/confinement}
\Longrightarrow \text{master equation}
\\
&\Longrightarrow \text{spectral-resolvent-semigroup-fusion synthesis}
\\
&\Longrightarrow \text{return readout of the same structure.}
\end{aligned}
\]
\end{remark}


\subsection*{Phase IV-P: Observable Back-Projection and Kernel Reconstruction}
\addcontentsline{toc}{subsection}{Phase IV-P: Observable Back-Projection and Kernel Reconstruction}

\begin{lemma}[Admissible back-projection lemma]\label{lem:admissible_back_projection}
If a late observable synthesis is admissible in the sense of Phases IV-J through IV-O, then it admits a back-projection into the grouped architecture such that the reconstructed object remains compatible with the same closure, transport, and non-drift constraints that governed the forward readout.
\end{lemma}
\begin{proof}
By Phase IV-N, the unified observable synthesis is not an arbitrary superposition of readout data, but a compatibility-preserving synthesis of spectral, resolvent, semigroup, and fusion constraints. Phase IV-O then shows that this synthesis remains a governed readout of the same master-equation sector. Hence any admissible late synthesis can be projected back through the grouped architecture without leaving the admissible class: what is reconstructed is not a new source, but the same structural body viewed through the inverse admissibility constraints of the readout chain.
\end{proof}

\begin{proposition}[Kernel reconstruction proposition]\label{prop:kernel_reconstruction}
Admissible back-projection reconstructs not a second kernel, but a readout-identified representative of the original generative kernel. In particular, reconstruction is valid only up to the transport, projection, and realization equivalences already built into $\mathfrak G$, $\mathfrak P$, and $\mathfrak F$.
\end{proposition}
\begin{proof}
The grouped transport architecture never licenses ontology drift; it only transports one admissible body across grouped layers. Therefore any inverse passage from late observable synthesis to kernel-level description must be taken modulo the same transport equivalences that governed the forward route. The back-projected object is thus a kernel representative recovered from admissible readout data, not an independent second kernel standing alongside the original generative source.
\end{proof}

\begin{theorem}[Observable-to-kernel closure theorem]\label{thm:observable_to_kernel_closure}
The late unfolding program closes in both directions: every admissible generative structure yields a governed late observable synthesis, and every admissible late observable synthesis yields a back-projected kernel representative of that same structure. Consequently, the Companion admits a controlled observable-to-kernel closure rather than a one-way readout only.
\end{theorem}
\begin{proof}
The forward direction is the content of the unfolding chain from the generative core through the master equation to the unified observable synthesis. By Lemma~\ref{lem:admissible_back_projection}, an admissible late synthesis can be projected back into the grouped architecture while preserving closure, transport, and non-drift. Proposition~\ref{prop:kernel_reconstruction} then shows that the reconstructed object is to be read as a representative of the original generative kernel, not as a second ontological origin. Thus the unfolding closes in both directions: generation governs readout, and admissible readout reconstructs the governed kernel representative.
\end{proof}

\begin{corollary}[Bidirectional closure chain]\label{cor:bidirectional_closure_chain}
The late-stage closure may therefore be compressed as
\[
\begin{aligned}
&\text{generative core}
\Longrightarrow \text{master-equation-governed readout}
\Longrightarrow \text{unified observable synthesis}\\
&\Longrightarrow \text{admissible back-projection}
\Longrightarrow \text{kernel representative of the same structure.}
\end{aligned}
\]
\end{corollary}


\subsection*{Phase IV-Q: Global Unfolding Synthesis and Architectural Compression}
\addcontentsline{toc}{subsection}{Phase IV-Q: Global Unfolding Synthesis and Architectural Compression}

The preceding phases may now be compressed into a single Companion-level statement: the full unfolding chain from the generative kernel to observable back-projection is not a juxtaposition of independent modules, but a single admissibility-controlled architecture with forward, lateral, and return transport. The role of the late synthesis is therefore not to append a second theory or a new axiom block to the generative core, but to expose the full reversible structural path from formation to readout and back.

\begin{proposition}[Global unfolding synthesis]\label{prop:global_unfolding_synthesis}
The Companion unfolding program may be read as one admissibility-preserving architecture
\[
\begin{aligned}
\text{generative formation}
&\Longrightarrow \text{HC/closure/confinement forcing}\\
&\Longrightarrow \text{cascade and readout admissibility}\\
&\Longrightarrow \text{semantic/tensor transport}\\
&\Longrightarrow \text{grouped-layer realization}\\
&\Longrightarrow \text{probe and field realization}\\
&\Longrightarrow \text{physical interface control}\\
&\Longrightarrow \text{spectral/resolvent/semigroup/fusion synthesis}\\
&\Longrightarrow \text{observable back-projection}\\
&\Longrightarrow \text{kernel reconstruction of the same structure}.
\end{aligned}
\]
In particular, the later readout layers do not introduce a second ontological source, but provide a controlled return map into the same structural kernel whose forcing chain generated the admissible dynamics in the first place.
\end{proposition}

\begin{proof}
Phase~I establishes the generative forcing chain from one-body discipline to HC, closure, and higher confinement. Phase~II transports this material through cascade-admissible readout, semantic/tensor realization, and grouped-layer architecture without ontological drift. Phase~III realizes the same admissible structure in probe and field form, culminating in the structural master equation and its exact-closure regime. Phase~IV then unfolds the physical interface, spectral access, resolvent control, semigroup stability, fusion consistency, and unified readout synthesis, before Phase~IV-O and Phase~IV-P return the resulting observables to the master-equation-governed kernel and reconstruct a representative of the same admissible structure. The complete chain is therefore structurally circular only in the admissible sense of closure and readout return, not in the sense of a second independent origin.\qedhere
\end{proof}

\begin{corollary}[Companion architectural compression rule]
Any later compression of the Companion is admissible only if it preserves explicit access to each of the following four structural roles:
\[
\text{formation},\qquad
\text{transport},\qquad
\text{realization/readout},\qquad
\text{return reconstruction}.
\]
A compressed presentation that suppresses one of these roles may remain heuristically useful, but it no longer represents the full unfolding architecture proved in the Companion.
\end{corollary}

\paragraph{Interpretive role of the global synthesis.}
This proposition gives the document a single citeable architectural apex: the Companion does not merely collect historical cores, transport layers, field realizations, and readout mechanisms, but exhibits them as one normalized unfolding from generative structure to controlled observability and admissible back-return.



\subsection*{Phase IV-R: Proof Compression and Structural Admissibility Testing}
\addcontentsline{toc}{subsection}{Phase IV-R: Proof Compression and Structural Admissibility Testing}

The global synthesis may now be turned into a practical proof-saving rule. Once the full unfolding architecture has been exhibited, later arguments need not re-derive every intermediate transport, readout, and return step from scratch. Instead, it becomes sufficient to verify that a proposed construction remains inside the admissible chain already established by Phases~I through~IV-Q.

\begin{lemma}[Structural admissibility test]\label{lem:structural_admissibility_test}
Let $X$ be a proposed extension, readout, or reformulation of the Companion architecture. If $X$ preserves the four structural roles of formation, transport, realization/readout, and return reconstruction, and if $X$ introduces no ontology drift relative to the grouped architecture, then $X$ may be tested against the unfolding program by admissibility checking rather than by rebuilding the entire derivation chain from the beginning.
\end{lemma}
\begin{proof}
By Corollary~[Companion architectural compression rule], any admissible compression or reformulation must preserve access to the four defining structural roles. Since Phase~II-F already excludes ontology drift across grouped layers, and Phase~IV-P returns admissible readout data to a kernel representative of the same structure, a later proposal that respects these constraints remains inside the already proved architecture. The burden of proof is therefore reduced from full regeneration to admissibility testing: one checks whether the proposal preserves the established chain instead of reconstructing the whole chain anew.
\end{proof}

\begin{proposition}[Proof-compression proposition]\label{prop:proof_compression_proposition}
The Companion permits controlled proof compression: whenever a later statement is structurally subordinate to an already proved admissible chain, it may be justified by citing preservation of that chain together with the specific additional admissibility condition that differentiates the later statement.
\end{proposition}
\begin{proof}
Phases~I through~IV-Q provide the explicit full chain from formation to observable back-return. If a later statement does not alter this chain, but only specializes one of its stages---for example by restricting admissible spectral modes, selecting a probe sector, or refining an operator gate---then the full derivation need not be repeated. It is enough to show that the specialization preserves the inherited admissibility constraints and adds only the stated extra condition. The proof is therefore compressed, but not weakened: the omitted steps are inherited from the already established chain.
\end{proof}

\begin{theorem}[Meta-theorem of admissible structural compression]\label{thm:meta_admissible_structural_compression}
Any later expansion of the Companion is mathematically legitimate in compressed form if and only if it can be located inside the established unfolding architecture without breaking formation, transport, realization/readout, or return reconstruction. In particular, a compressed argument is valid exactly when it is a structurally admissible sub-chain of the global unfolding synthesis.
\end{theorem}
\begin{proof}
Necessity follows from the architectural compression rule: if one of the four structural roles is broken, the proposal no longer belongs to the full unfolding architecture and cannot inherit its proof force. Sufficiency follows from Lemma~\ref{lem:structural_admissibility_test} and Proposition~\ref{prop:proof_compression_proposition}: once a proposal is shown to preserve the admissible chain and differs only by an explicit additional admissibility condition, the rest of the proof is inherited from the already established unfolding program. Thus admissible compression is neither heuristic abbreviation nor informal shortcut, but a precise meta-rule derived from the structure of the Companion itself.
\end{proof}

\begin{corollary}[Late meta-compression chain]\label{cor:late_meta_compression_chain}
The full late-stage architecture may therefore be compressed as
\[
\begin{aligned}
&\text{global unfolding synthesis}
\Longrightarrow \text{admissibility testing of later proposals}\\
&\Longrightarrow \text{controlled proof compression}
\Longrightarrow \text{valid structural sub-chain expansions.}
\end{aligned}
\]
\end{corollary}

\paragraph{Interpretive role of the meta-compression theorem.}
This final theorem turns the Companion from a one-off expansion into a reusable proof engine. Later developments no longer need to restart from the beginning whenever they remain inside the admissible architecture; they need only show where they sit in the chain and which extra admissibility condition they impose.



% --- inserted late main-text phases IV-S and IV-T ---

%% ============================================================
%%  s_op1_op4_hopf_ckm.tex
%%  STATUS: PROOF MODE — Erster strukturierter Beweisentwurf
%%
%%  OP 1: β_sym = 1/√3 aus Hopf-Fixpunkt-Geometrie
%%  OP 4: CKM/PMNS aus G₂ Euler-Winkeln (Faktor ≈ 2 schließen)
%%
%%  Marc Brendecke, Bochum, März 2026
%%  Drop-In für quantum_sphaera v2.x (nach Section 11)
%%
%%  SCHLÜSSELEINSICHT: OP 1 und OP 4 sind dasselbe Problem.
%%  Der Faktor ≈ 2 in den Mischungswinkeln (OP 4) ist dieselbe
%%  5.1%-Lücke wie in β_sym (OP 1): beides ist die
%%  Kontinuumslimit-Korrektur der diskreten Primgitter-Geometrie.
%% ============================================================

\subsection*{Phase IV-S: Hopf Fixed-Point Geometry and CKM/PMNS Residuals}
\addcontentsline{toc}{subsection}{Phase IV-S: Hopf Fixed-Point Geometry and CKM/PMNS Residuals}
\label{sec:hopf_ckm}

\begin{remark}[Transition from Phase IV-R to the Hopf/CKM block]\label{rem:transition_r_to_s}
Phase~IV-R turned the late synthesis into a reusable admissibility and proof-compression rule. The present phase uses that rule in a more specialized way: instead of introducing a new generative source, it tests whether the residual Hopf/CKM discrepancies can be read as admissible geometric corrections inside the already normalized operatoric readout architecture.
\end{remark}

